\section{Virtualização de Funções de Rede (NFV)}

Com o crescente número de funções presentes nas redes, a dependência de equipamentos específicos pode aumentar a complexidade do gerenciamento e do fornecimento de novos serviços. Nesse contexto, a Virtualização de Funções de Rede (NFV, Network Function Virtualization) permite desacoplar determinadas funções de rede do hardware dedicado, possibilitando sua implementação em software sobre uma infraestrutura de propósito geral \cite{yi2018comprehensive}.

As funções de rede implementadas de forma virtualizada são denominadas Funções Virtuais de Rede (VNFs, Virtual Network Functions) e podem desempenhar funcionalidades equivalentes às tradicionalmente executadas por equipamentos físicos dedicados. Ao reduzir a dependência entre essas funções e um hardware específico, a NFV proporciona maior flexibilidade para a implantação e o gerenciamento de serviços, permitindo que diferentes funções sejam disponibilizadas de acordo com as necessidades da rede \cite{yi2018comprehensive}.

A possibilidade de implementar funções de rede de forma virtualizada contribui para a criação de serviços com diferentes configurações sobre uma infraestrutura compartilhada. No contexto do fatiamento de redes, a NFV permite que diferentes funções sejam instanciadas e organizadas de acordo com os requisitos de cada fatia, contribuindo para sua flexibilidade e adaptação. Dessa forma, em conjunto com abordagens como NV e SDN, a NFV fornece mecanismos importantes para a implementação e o gerenciamento de fatias de rede \cite{gomes2020strategies}.